\documentclass[10pt]{article}
\usepackage[paperwidth=216mm, paperheight=279mm, left=64mm, right=20mm, top=22mm, bottom=20mm, headheight=14pt, headsep=6mm]{geometry}
\usepackage{fontspec}
\usepackage{amsmath,amssymb}
\usepackage{xcolor}
\usepackage{fancyhdr}
\usepackage{parskip}

\definecolor{stewartblue}{RGB}{0, 130, 145}

\pagestyle{fancy}
\fancyhf{}
\renewcommand{\headrulewidth}{0pt}
\fancyhead[L]{\small\sffamily\textbf{MỤC 3.7}\quad Tốc độ thay đổi trong Khoa học Tự nhiên và Xã hội}
\fancyhead[R]{\small\sffamily\textbf{233}}

\setlength{\parindent}{1.5em}
\setlength{\parskip}{0.35em plus 0.1em minus 0.1em}

\begin{document}

\noindent
sang micromét ($1\text{ cm} = 10.000\ \mu\text{m}$). Khi đó bán kính của động mạch là $80\ \mu\text{m}$. Vận tốc tại trục tâm là $11.850\ \mu\text{m/s}$, giảm xuống còn $11.110\ \mu\text{m/s}$ ở khoảng cách $r = 20\ \mu\text{m}$. Việc $dv/dr = -74\ (\mu\text{m/s})/\mu\text{m}$ có nghĩa là, khi $r = 20\ \mu\text{m}$, vận tốc đang giảm với tốc độ khoảng $74\ \mu\text{m/s}$ cho mỗi micromét khi chúng ta di chuyển ra xa tâm.\hfill{\color{stewartblue}\rule{1.1ex}{1.1ex}}

\vspace{1.2em}

\noindent{\color{stewartblue}\rule[0.05ex]{1.15ex}{1.15ex}}\hspace{0.5em}{\sffamily\bfseries\large Kinh tế học}

\vspace{0.8em}

\noindent{\sffamily\bfseries\color{stewartblue}VÍ DỤ 8}\quad Giả sử $C(x)$ là tổng chi phí mà một công ty phải chịu khi sản xuất $x$ đơn vị của một loại hàng hóa nào đó. Hàm $C$ được gọi là \textbf{hàm chi phí}. Nếu số lượng sản phẩm sản xuất tăng từ $x_1$ lên $x_2$, thì chi phí phát sinh thêm là $\Delta C = C(x_2) - C(x_1)$, và tốc độ thay đổi trung bình của chi phí là
\[
\frac{\Delta C}{\Delta x} = \frac{C(x_2) - C(x_1)}{x_2 - x_1} = \frac{C(x_1 + \Delta x) - C(x_1)}{\Delta x}
\]

Giới hạn của đại lượng này khi $\Delta x \to 0$, tức là tốc độ thay đổi tức thời của chi phí theo số lượng sản phẩm được sản xuất, được các nhà kinh tế học gọi là \textbf{chi phí biên}:
\[
\text{chi phí biên} = \lim_{\Delta x \to 0} \frac{\Delta C}{\Delta x} = \frac{dC}{dx}
\]

[Vì $x$ thường chỉ nhận các giá trị nguyên, việc cho $\Delta x$ tiến dần về $0$ có thể không mang ý nghĩa thực tế theo nghĩa đen, nhưng chúng ta luôn có thể thay thế $C(x)$ bằng một hàm xấp xỉ trơn như trong Ví dụ 6.]

Lấy $\Delta x = 1$ và $n$ lớn (sao cho $\Delta x$ là nhỏ so với $n$), ta có
\[
C'(n) \approx C(n + 1) - C(n)
\]

Như vậy, chi phí biên của việc sản xuất $n$ đơn vị xấp xỉ bằng chi phí sản xuất thêm một đơn vị nữa [đơn vị thứ $(n + 1)$].

Người ta thường biểu diễn một hàm tổng chi phí bằng một đa thức
\[
C(x) = a + bx + cx^2 + dx^3
\]
trong đó $a$ đại diện cho chi phí gián tiếp cố định (tiền thuê mặt bằng, sưởi ấm, bảo trì) và các số hạng khác đại diện cho chi phí nguyên vật liệu, nhân công, vân vân. (Chi phí nguyên vật liệu có thể tỷ lệ thuận với $x$, nhưng chi phí nhân công có thể phụ thuộc một phần vào các lũy thừa bậc cao hơn của $x$ do chi phí làm thêm giờ và sự kém hiệu quả phát sinh trong các vận hành quy mô lớn.)

Chẳng hạn, giả sử một công ty đã ước tính rằng chi phí (tính bằng đô la) để sản xuất $x$ sản phẩm là
\[
C(x) = 10.000 + 5x + 0{,}01x^2
\]
Khi đó hàm chi phí biên là
\[
C'(x) = 5 + 0{,}02x
\]
Chi phí biên ở mức sản xuất 500 sản phẩm là
\[
C'(500) = 5 + 0{,}02(500) = \$15/\text{sản phẩm}
\]

\vfill

\begin{flushleft}
{\fontsize{5.5}{6.5}\selectfont\color{gray!70}
Copyright 2021 Cengage Learning. All Rights Reserved. May not be copied, scanned, or duplicated, in whole or in part. WCN 02-200-203
\par}
\end{flushleft}

\end{document}